\documentclass[aspectratio=169, 12pt]{beamer}
% aspectratio: 169 для 16:9, 43 для 4:3

% ============================================
% ТЕМА ОФОРМЛЕНИЯ
% ============================================
\usetheme{Madrid}
% Другие темы: Berlin, Copenhagen, Warsaw, Singapore, CambridgeUS

\usecolortheme{default}
% Другие цветовые схемы: dolphin, beaver, crane, seahorse, wolverine

% Настройка цветов
\definecolor{myblue}{RGB}{0, 51, 102}
\definecolor{myred}{RGB}{153, 0, 0}
\definecolor{mygreen}{RGB}{0, 102, 51}

\setbeamercolor{structure}{fg=myblue}
\setbeamercolor{title}{fg=white, bg=myblue}
\setbeamercolor{frametitle}{fg=white, bg=myblue}
\setbeamercolor{block title}{fg=white, bg=myblue}
\setbeamercolor{block body}{bg=myblue!10}
\setbeamercolor{block title example}{fg=white, bg=mygreen}
\setbeamercolor{block body example}{bg=mygreen!10}
\setbeamercolor{block title alerted}{fg=white, bg=myred}
\setbeamercolor{block body alerted}{bg=myred!10}

% ============================================
% КОДИРОВКА И ЯЗЫК
% ============================================
\usepackage[utf8]{inputenc}
\usepackage[T2A]{fontenc}
\usepackage[russian]{babel}

% ============================================
% ДОПОЛНИТЕЛЬНЫЕ ПАКЕТЫ
% ============================================

% Графика
\usepackage{graphicx}
\graphicspath{{images/}}

% Таблицы
\usepackage{booktabs}
\usepackage{array}

% Математика
\usepackage{amsmath, amssymb}

% Листинги кода
\usepackage{listings}
\lstset{
  basicstyle=\ttfamily\footnotesize,
  breaklines=true,
  frame=single,
  backgroundcolor=\color{gray!10},
  commentstyle=\color{green!50!black},
  keywordstyle=\color{blue},
  stringstyle=\color{red!50!black},
  showstringspaces=false,
  tabsize=2,
  language=TeX
}

% Гиперссылки
\usepackage{hyperref}
\hypersetup{
  colorlinks=true,
  linkcolor=myblue,
  urlcolor=mygreen,
  citecolor=myred
}

% Иконки и символы
\usepackage{fontawesome5}

% Многоколоночная верстка
\usepackage{multicol}

% ============================================
% БИБЛИОГРАФИЯ
% ============================================
\usepackage[
  backend=biber,
  style=authoryear,
  maxcitenames=2,
  maxbibnames=3
]{biblatex}

\addbibresource{references.bib}

% Настройка библиографии для презентации
\setbeamertemplate{bibliography item}{\insertbiblabel}
\renewcommand*{\bibfont}{\footnotesize}

% ============================================
% НАСТРОЙКА ВНЕШНЕГО ВИДА
% ============================================

% Убрать навигационные символы
\setbeamertemplate{navigation symbols}{}

% Номера слайдов
\setbeamertemplate{footline}[frame number]

% Стиль маркеров списков
\setbeamertemplate{itemize items}[circle]
\setbeamertemplate{enumerate items}[default]

% Оглавление в начале каждой секции
\AtBeginSection[]
{
  \begin{frame}
    \frametitle{План презентации}
    \tableofcontents[currentsection]
  \end{frame}
}

% ============================================
% ИНФОРМАЦИЯ О ПРЕЗЕНТАЦИИ
% ============================================
\title[Библиография в LaTeX]{Работа с библиографией в LaTeX}
\subtitle{Использование BibTeX и biblatex}

\author[Викторов Е.И.]{
  Викторов Егор \\
  \small Группа НПМмд02-24
}

\institute[Университет]{
  Федеральное государственное бюджетное образовательное учреждение \\
  высшего образования \\
  \textbf{«РУДН»} \\
  \vspace{0.3cm}
  Факультет физико-математических и естественных наук \\
  Кафедра Теории Вероятностей и Кибербезопасности
}

\date{\today}

% Логотип университета (раскомментируйте и укажите путь к файлу)
% \logo{\includegraphics[height=1cm]{logo.png}}

% ============================================
% НАЧАЛО ПРЕЗЕНТАЦИИ
% ============================================
\begin{document}

% ============================================
% ТИТУЛЬНЫЙ СЛАЙД
% ============================================
\begin{frame}
  \titlepage
\end{frame}
% ============================================
% СОДЕРЖАНИЕ
% ============================================
\begin{frame}
  \frametitle{План презентации}
  \tableofcontents
\end{frame}

% ============================================
% ВВЕДЕНИЕ
% ============================================
\section{Введение}

\begin{frame}{Актуальность темы}
  \begin{columns}[T]
    \begin{column}{0.5\textwidth}
      \textbf{Проблемы ручного оформления:}
      \begin{itemize}
        \item \faTimesCircle\ Трудоемкость
        \item \faTimesCircle\ Высокая вероятность ошибок
        \item \faTimesCircle\ Сложность изменения стиля
        \item \faTimesCircle\ Неединообразие оформления
      \end{itemize}
    \end{column}
    
    \begin{column}{0.5\textwidth}
      \textbf{Преимущества автоматизации:}
      \begin{itemize}
        \item \faCheckCircle\ Экономия времени
        \item \faCheckCircle\ Минимум ошибок
        \item \faCheckCircle\ Легкая смена стиля
        \item \faCheckCircle\ Единый формат
      \end{itemize}
    \end{column}
  \end{columns}
  
  \vspace{1cm}
  
  \begin{alertblock}{Вывод}
    Автоматизация управления библиографией — необходимость для современных научных публикаций
  \end{alertblock}
\end{frame}

\begin{frame}{Цель и задачи работы}
  \begin{block}{Цель работы}
    Изучение методов работы с библиографическими ссылками в системе LaTeX и освоение инструментов BibTeX и biblatex
  \end{block}
  
  \vspace{0.5cm}
  
  \textbf{Задачи работы:}
  \begin{enumerate}
    \item Изучить структуру файлов библиографии формата \texttt{.bib}
    \item Освоить различные типы библиографических записей
    \item Научиться использовать команды цитирования
    \item Изучить различные стили оформления библиографии
    \item Настроить автоматическую компиляцию документа
    \item Создать документ с корректно оформленными ссылками
  \end{enumerate}
\end{frame}

% ============================================
% ТЕОРЕТИЧЕСКАЯ ЧАСТЬ
% ============================================
\section{Теоретическая часть}

\begin{frame}{История развития}
  \begin{center}
    \begin{tikzpicture}[scale=0.9]
      % Timeline
      \draw[thick, ->] (0,0) -- (12,0);
      
      % Events
      \draw[thick] (2,0) -- (2,0.3);
      \node[above] at (2,0.3) {\small 1978};
      \node[below, text width=2cm, align=center] at (2,-0.5) {\tiny TeX \\ (Knuth)};
      
      \draw[thick] (5,0) -- (5,0.3);
      \node[above] at (5,0.3) {\small 1985};
      \node[below, text width=2cm, align=center] at (5,-0.5) {\tiny LaTeX \\ (Lamport)};
      
      \draw[thick] (8,0) -- (8,0.3);
      \node[above] at (8,0.3) {\small 1988};
      \node[below, text width=2cm, align=center] at (8,-0.5) {\tiny BibTeX \\ (Patashnik)};
      
      \draw[thick] (11,0) -- (11,0.3);
      \node[above] at (11,0.3) {\small 2006};
      \node[below, text width=2cm, align=center] at (11,-0.5) {\tiny biblatex \\ (Lehman)};
    \end{tikzpicture}
  \end{center}
  
  \vspace{0.5cm}
  
  \begin{itemize}
    \item \textbf{1978} — Дональд Кнут создает TeX
    \item \textbf{1985} — Лесли Лампорт разрабатывает LaTeX
    \item \textbf{1988} — Орен Паташник представляет BibTeX
    \item \textbf{2006} — Филипп Леман создает biblatex
  \end{itemize}
\end{frame}

\begin{frame}{BibTeX vs biblatex}
  \begin{columns}[T]
    \begin{column}{0.48\textwidth}
      \begin{block}{BibTeX}
        \textbf{Преимущества:}
        \begin{itemize}
          \item Простота
          \item Стабильность
          \item Широкая поддержка
        \end{itemize}
        
        \vspace{0.3cm}
        
        \textbf{Недостатки:}
        \begin{itemize}
          \item Ограниченная Unicode поддержка
          \item Сложность настройки
          \item Устаревший код
        \end{itemize}
      \end{block}
    \end{column}
    
    \begin{column}{0.48\textwidth}
      \begin{exampleblock}{biblatex}
        \textbf{Преимущества:}
        \begin{itemize}
          \item Полная Unicode поддержка
          \item Гибкая настройка
          \item Современный подход
          \item Backend Biber
        \end{itemize}
        
        \vspace{0.3cm}
        
        \textbf{Недостатки:}
        \begin{itemize}
          \item Более сложный
          \item Требует изучения
        \end{itemize}
      \end{exampleblock}
    \end{column}
  \end{columns}
\end{frame}

\begin{frame}[fragile]{Структура файла библиографии}
  \textbf{Общая структура записи:}
  
  \begin{lstlisting}
@ТИП_ИСТОЧНИКА{ключ_цитирования,
  поле1 = {значение1},
  поле2 = {значение2},
  поле3 = {значение3}
}
  \end{lstlisting}
  
  \vspace{0.5cm}
  
  \textbf{Пример записи книги:}
  
  \begin{lstlisting}
@book{knuth1984,
  author    = {Donald E. Knuth},
  title     = {The TeXbook},
  publisher = {Addison-Wesley},
  year      = {1984},
  isbn      = {0-201-13447-0}
}
  \end{lstlisting}
\end{frame}

\begin{frame}{Типы библиографических записей}
  \begin{columns}[T]
    \begin{column}{0.48\textwidth}
      \textbf{Печатные издания:}
      \begin{itemize}
        \item \texttt{@book} — книга
        \item \texttt{@article} — статья
        \item \texttt{@inproceedings} — конференция
        \item \texttt{@phdthesis} — диссертация
        \item \texttt{@manual} — руководство
      \end{itemize}
    \end{column}
    
    \begin{column}{0.48\textwidth}
      \textbf{Электронные ресурсы:}
      \begin{itemize}
        \item \texttt{@online} — веб-сайт
        \item \texttt{@techreport} — тех. отчет
        \item \texttt{@misc} — прочее
        \item \texttt{@unpublished} — неопубликованное
      \end{itemize}
    \end{column}
  \end{columns}
  
  \vspace{1cm}
  
  \begin{exampleblock}{Совет}
    Выбирайте наиболее подходящий тип для каждого источника — это влияет на форматирование в списке литературы
  \end{exampleblock}
\end{frame}

\begin{frame}{Команды цитирования}
  \begin{table}
    \centering
    \small
    \begin{tabular}{@{}lll@{}}
      \toprule
      \textbf{Команда} & \textbf{Результат} & \textbf{Применение} \\
      \midrule
      \verb|\cite{key}| & [1] & Базовое \\
      \verb|\parencite{key}| & (Knuth, 1984) & В скобках \\
      \verb|\textcite{key}| & Knuth (1984) & В тексте \\
      \verb|\citeauthor{key}| & Knuth & Только автор \\
      \verb|\citeyear{key}| & 1984 & Только год \\
      \verb|\footcite{key}| & \textsuperscript{1} & В сноске \\
      \bottomrule
    \end{tabular}
  \end{table}
  
  \vspace{0.5cm}
  
  \begin{alertblock}{Важно!}
    Выбор команды зависит от стиля цитирования и контекста использования
  \end{alertblock}
\end{frame}

% ============================================
% ПРАКТИЧЕСКАЯ ЧАСТЬ
% ============================================
\section{Практическая часть}

\begin{frame}[fragile]{Настройка документа}
  \textbf{Подключение пакета biblatex:}
  
  \begin{lstlisting}
\usepackage[
  backend=biber,        % Используем Biber
  style=gost-numeric,   % Стиль ГОСТ
  sorting=none,         % Без сортировки
  maxbibnames=99,       % Все авторы
  doi=true,             % Показывать DOI
  isbn=true,            % Показывать ISBN
  url=true              % Показывать URL
]{biblatex}

\addbibresource{references.bib}
  \end{lstlisting}
  
  \vspace{0.3cm}
  
  \textbf{Вывод списка литературы:}
  
  \begin{lstlisting}
\printbibliography[title={Список литературы}]
  \end{lstlisting}
\end{frame}

\begin{frame}{Процесс компиляции}
  \begin{center}
    \begin{tikzpicture}[
      node distance=1.5cm,
      box/.style={rectangle, draw, fill=blue!20, text width=3cm, align=center, minimum height=1cm},
      arrow/.style={->, thick}
    ]
      \node[box] (step1) {1. pdflatex};
      \node[box, below of=step1] (step2) {2. biber};
      \node[box, below of=step2] (step3) {3. pdflatex};
      \node[box, below of=step3] (step4) {4. pdflatex};
      
      \draw[arrow] (step1) -- (step2);
      \draw[arrow] (step2) -- (step3);
      \draw[arrow] (step3) -- (step4);
      \node[right=0.5cm of step1, text width=4cm] {Создание .aux файла};
      \node[right=0.5cm of step2, text width=4cm] {Обработка библиографии};
      \node[right=0.5cm of step3, text width=4cm] {Включение библиографии};
      \node[right=0.5cm of step4, text width=4cm] {Обновление ссылок};
    \end{tikzpicture}
  \end{center}
\end{frame}

\begin{frame}[fragile]{Автоматизация с latexmk}
  \textbf{Создание файла .latexmkrc:}
  
  \begin{lstlisting}[language=bash]
$pdf_mode = 1;
$biber = 'biber %O %S';
$pdflatex = 'pdflatex -interaction=nonstopmode';
  \end{lstlisting}
  
  \vspace{0.5cm}
  
  \textbf{Команды компиляции:}
  
  \begin{lstlisting}[language=bash]
# Компиляция
latexmk -pdf document.tex

# Очистка временных файлов
latexmk -c

# Полная очистка
latexmk -C
  \end{lstlisting}
\end{frame}

\begin{frame}{Примеры цитирования}
  \begin{exampleblock}{Цитирование книги}
    Классическая книга по TeX \parencite{knuth1984} является фундаментальным трудом.
  \end{exampleblock}
  
  \begin{exampleblock}{Цитирование статьи}
    Историческая статья \textcite{dijkstra1968} оказала значительное влияние.
  \end{exampleblock}
  
  \begin{exampleblock}{Множественное цитирование}
    Основы LaTeX описаны во многих источниках \parencite{knuth1984, lamport1994, mittelbach2004}.
  \end{exampleblock}
  
  \begin{exampleblock}{Онлайн-ресурс}
    Официальная документация доступна на сайте \parencite{latex-project}.
  \end{exampleblock}
\end{frame}

\begin{frame}{Стили оформления библиографии}
  \begin{columns}[T]
    \begin{column}{0.48\textwidth}
      \textbf{Числовые стили:}
      \begin{itemize}
        \item \texttt{numeric} — [1], [2], [3]
        \item \texttt{ieee} — стиль IEEE
        \item \texttt{gost-numeric} — ГОСТ
      \end{itemize}
      
      \vspace{0.5cm}
      
      \textbf{Алфавитные стили:}
      \begin{itemize}
        \item \texttt{alphabetic} — [Knu84]
        \item \texttt{authoryear} — (Knuth, 1984)
        \item \texttt{gost-authoryear} — ГОСТ
      \end{itemize}
    \end{column}
    
    \begin{column}{0.48\textwidth}
      \textbf{Специализированные:}
      \begin{itemize}
        \item \texttt{apa} — APA стиль
        \item \texttt{mla} — MLA стиль
        \item \texttt{chicago} — Чикаго
        \item \texttt{nature} — Nature
        \item \texttt{science} — Science
      \end{itemize}
    \end{column}
  \end{columns}
  
  \vspace{0.5cm}
  
  \begin{alertblock}{Примечание}
    Стиль легко изменить, изменив один параметр в настройках biblatex
  \end{alertblock}
\end{frame}

% ============================================
% АНАЛИЗ РЕЗУЛЬТАТОВ
% ============================================
\section{Анализ результатов}

\begin{frame}{Сравнение методов оформления}
  \begin{table}
    \centering
    \footnotesize
    \begin{tabular}{@{}lcc@{}}
      \toprule
      \textbf{Критерий} & \textbf{BibTeX/biblatex} & \textbf{Ручное} \\
      \midrule
      Время на оформление & \textcolor{green}{\faArrowDown\ Низкое} & \textcolor{red}{\faArrowUp\ Высокое} \\
      Вероятность ошибок & \textcolor{green}{\faArrowDown\ Низкая} & \textcolor{red}{\faArrowUp\ Высокая} \\
      Единообразие & \textcolor{green}{\faCheck\ Да} & \textcolor{red}{\faTimes\ Требует внимания} \\
      Изменение стиля & \textcolor{green}{\faCheck\ Легко} & \textcolor{red}{\faTimes\ Сложно} \\
      Переиспользование & \textcolor{green}{\faCheck\ Легко} & \textcolor{red}{\faTimes\ Сложно} \\
      Автосортировка & \textcolor{green}{\faCheck\ Да} & \textcolor{red}{\faTimes\ Нет} \\
      \bottomrule
    \end{tabular}
  \end{table}
  
  \vspace{0.5cm}
  
  \begin{block}{Вывод}
    Автоматизация превосходит ручное оформление по всем критериям
  \end{block}
\end{frame}

\begin{frame}{Преимущества biblatex}
  \begin{columns}[c]
    \begin{column}{0.48\textwidth}
      \begin{block}{Технические}
        \begin{itemize}
          \item \faCode\ Полная Unicode поддержка
          \item \faCogs\ Гибкая настройка
          \item \faDatabase\ Мощный backend Biber
          \item \faLanguage\ Локализация
        \end{itemize}
      \end{block}
    \end{column}
    
    \begin{column}{0.48\textwidth}
      \begin{block}{Практические}
        \begin{itemize}
          \item \faClock\ Экономия времени
          \item \faCheckDouble\ Минимум ошибок
          \item \faSync\ Переиспользование
          \item \faPalette\ Единый стиль
        \end{itemize}
      \end{block}
    \end{column}
  \end{columns}
  
  \vspace{1cm}
  
  \begin{center}
    \Large
    \textcolor{mygreen}{\faThumbsUp} Рекомендуется для всех новых проектов!
  \end{center}
\end{frame}

\begin{frame}{Рекомендации по использованию}
  \begin{enumerate}
    \item \textbf{Выбор инструмента}
    \begin{itemize}
      \item Используйте biblatex для новых проектов
      \item Выбирайте backend=biber для Unicode
    \end{itemize}
    
    \item \textbf{Организация файлов}
    \begin{itemize}
      \item Храните библиографию в отдельном .bib файле
      \item Используйте осмысленные ключи цитирования
    \end{itemize}
    
    \item \textbf{Инструменты}
    \begin{itemize}
      \item JabRef или Zotero для управления базой
      \item latexmk для автоматической компиляции
    \end{itemize}
    
    \item \textbf{Безопасность}
    \begin{itemize}
      \item Регулярно делайте резервные копии
      \item Используйте систему контроля версий (Git)
    \end{itemize}
  \end{enumerate}
\end{frame}

\begin{frame}{Типичные ошибки и их решение}
  \begin{alertblock}{Ошибка 1: Библиография не появляется}
    \textbf{Решение:} Проверьте последовательность компиляции (pdflatex → biber → pdflatex → pdflatex)
  \end{alertblock}
  
  \begin{alertblock}{Ошибка 2: Неправильная кодировка}
    \textbf{Решение:} Убедитесь, что .bib файл сохранен в UTF-8
  \end{alertblock}
  
  \begin{alertblock}{Ошибка 3: Источник не цитируется}
    \textbf{Решение:} Проверьте правильность ключа цитирования и наличие записи в .bib
  \end{alertblock}
  
  \begin{alertblock}{Ошибка 4: Неправильный формат}
    \textbf{Решение:} Проверьте синтаксис записи в .bib файле (запятые, скобки)
  \end{alertblock}
\end{frame}

% ============================================
% ЗАКЛЮЧЕНИЕ
% ============================================
\section{Заключение}

\begin{frame}{Результаты работы}
  \textbf{Выполнено:}
  
  \begin{itemize}
    \item[\faCheckCircle] Изучена структура файлов библиографии
    \item[\faCheckCircle] Освоены типы библиографических записей
    \item[\faCheckCircle] Изучены команды цитирования
    \item[\faCheckCircle] Рассмотрены различные стили оформления
    \item[\faCheckCircle] Настроена автоматическая компиляция
    \item[\faCheckCircle] Создан документ с корректными ссылками
  \end{itemize}
  
  \vspace{0.5cm}
  
  \begin{block}{Итог}
    Все поставленные задачи выполнены, цель работы достигнута
  \end{block}
\end{frame}

\begin{frame}{Практическая значимость}
  \textbf{Полученные навыки применимы для:}
  
  \begin{multicols}{2}
    \begin{itemize}
      \item Курсовых работ
      \item Дипломных проектов
      \item Научных статей
      \item Диссертаций
      \item Технических отчетов
      \item Книг и монографий
      \item Презентаций
      \item Учебных пособий
    \end{itemize}
  \end{multicols}
  
  \vspace{0.5cm}
  
  \begin{exampleblock}{Перспективы}
    Освоение biblatex открывает возможности для профессионального оформления любых научных и технических документов
  \end{exampleblock}
\end{frame}

\begin{frame}{Дальнейшее развитие}
  \textbf{Направления для изучения:}
  
  \begin{enumerate}
    \item \textbf{Продвинутые возможности biblatex}
    \begin{itemize}
      \item Создание собственных стилей
      \item Настройка полей и форматов
      \item Работа с мультиязычными источниками
    \end{itemize}
    
    \item \textbf{Интеграция с другими инструментами}
    \begin{itemize}
      \item Менеджеры библиографии (JabRef, Zotero)
      \item Системы контроля версий (Git)
      \item Онлайн-редакторы (Overleaf)
    \end{itemize}
    
    \item \textbf{Автоматизация рабочего процесса}
    \begin{itemize}
      \item Скрипты для обработки библиографии
      \item CI/CD для LaTeX документов
      \item Шаблоны и заготовки
    \end{itemize}
  \end{enumerate}
\end{frame}

% ============================================
% ПОЛЕЗНЫЕ РЕСУРСЫ
% ============================================
\section{Полезные ресурсы}

\begin{frame}{Официальная документация}
  \begin{block}{Основные ресурсы}
    \begin{itemize}
      \item \faGlobe\ The LaTeX Project: \url{https://www.latex-project.org/}
      \item \faBook\ CTAN: \url{https://www.ctan.org/}
      \item \faFileAlt\ Biblatex manual: \url{https://ctan.org/pkg/biblatex}
      \item \faCog\ Biber manual: \url{https://ctan.org/pkg/biber}
    \end{itemize}
  \end{block}
  
  \begin{block}{Обучающие материалы}
    \begin{itemize}
      \item \faGraduationCap\ Overleaf Learn: \url{https://www.overleaf.com/learn}
      \item \faQuestionCircle\ TeX Stack Exchange: \url{https://tex.stackexchange.com/}
      \item \faUsers\ LaTeX.ru: \url{https://www.latex.ru/}
    \end{itemize}
  \end{block}
\end{frame}

\begin{frame}{Инструменты}
  \begin{columns}[T]
    \begin{column}{0.48\textwidth}
      \begin{block}{Менеджеры библиографии}
        \begin{itemize}
          \item \faDatabase\ JabRef
          \item \faBookmark\ Zotero
          \item \faBook\ Mendeley
          \item \faPaperclip\ EndNote
        \end{itemize}
      \end{block}
      
      \begin{block}{Редакторы LaTeX}
        \begin{itemize}
          \item \faEdit\ TeXstudio
          \item \faCode\ TeXmaker
          \item \faCloud\ Overleaf
          \item \faTerminal\ VS Code + LaTeX Workshop
        \end{itemize}
      \end{block}
    \end{column}
    
    \begin{column}{0.48\textwidth}
      \begin{block}{Дистрибутивы TeX}
        \begin{itemize}
          \item \faWindows\ MiKTeX (Windows)
          \item \faLinux\ TeX Live (Linux)
          \item \faApple\ MacTeX (macOS)
        \end{itemize}
      \end{block}
      
      \begin{block}{Утилиты}
        \begin{itemize}
          \item \faSyncAlt\ latexmk
          \item \faFileCode\ arara
          \item \faGitAlt\ Git
          \item \faTerminal\ make
        \end{itemize}
      \end{block}
    \end{column}
  \end{columns}
\end{frame}

% ============================================
% ВОПРОСЫ
% ============================================
\begin{frame}[plain]
  \begin{center}
    \Huge \textcolor{myblue}{Спасибо за внимание!}
    
    \vspace{1cm}
    
    \Large \textcolor{mygreen}{Вопросы?}
    
    \vspace{2cm}
    
    \normalsize
    \begin{tabular}{rl}
     \end{tabular}
  \end{center}
\end{frame}

% ============================================
% СПИСОК ЛИТЕРАТУРЫ
% ============================================
\begin{frame}[allowframebreaks]{Список литературы}
  \printbibliography[heading=none]
\end{frame}

% ============================================
% ПРИЛОЖЕНИЯ
% ============================================
\appendix

\begin{frame}[fragile]{Приложение A: Пример .bib файла}
  \begin{lstlisting}[basicstyle=\ttfamily\tiny]
@book{knuth1984,
  author    = {Donald E. Knuth},
  title     = {The TeXbook},
  publisher = {Addison-Wesley Professional},
  year      = {1984},
  isbn      = {0-201-13447-0}
}

@article{dijkstra1968,
  author  = {Edsger W. Dijkstra},
  title   = {Go To Statement Considered Harmful},
  journal = {Communications of the ACM},
  year    = {1968},
  volume  = {11},
  number  = {3},
  pages   = {147--148}
}

@online{latex-project,
  author  = {{The LaTeX Project}},
  title   = {LaTeX – A document preparation system},
  year    = {2024},
  url     = {https://www.latex-project.org/}
}
  \end{lstlisting}
\end{frame}

\begin{frame}[fragile]{Приложение B: Команды компиляции}
  \textbf{Вариант 1: Ручная компиляция}
  \begin{lstlisting}[language=bash, basicstyle=\ttfamily\small]
pdflatex document.tex
biber document
pdflatex document.tex
pdflatex document.tex
  \end{lstlisting}
  
  \vspace{0.5cm}
  
  \textbf{Вариант 2: С использованием latexmk}
  \begin{lstlisting}[language=bash, basicstyle=\ttfamily\small]
latexmk -pdf -bibtex document.tex
  \end{lstlisting}
  
  \vspace{0.5cm}
  
  \textbf{Вариант 3: Makefile}
  \begin{lstlisting}[language=make, basicstyle=\ttfamily\small]
all:
 pdflatex document.tex
 biber document
 pdflatex document.tex
 pdflatex document.tex
  \end{lstlisting}
\end{frame}

\end{document}